% Generated from paper/full_draft. Edit the Markdown source or migration script.

\section{Conservative Baseline and Seed Stability Diagnostics}
\label{sec:12-appendix:conservative-baseline-and-seed-stability-diagnostics}

The conservative confidence-only context policy is the stable baseline for context compaction. The main text reports calibrated token-budget policies as the primary cost result; this supplement keeps the earlier confidence-only scale table and seed diagnostics. These intervals are engineering stability diagnostics, not strong inferential proof: each scale has only three observations. Tables~\ref{tab:s1} and~\ref{tab:s2} report quality and token stability.

\begin{table}[pos=tbp]
\centering
\small
\caption{Multi-seed retrieval-quality stability}
\label{tab:s1}
\begin{tabularx}{\textwidth}{>{\raggedright\arraybackslash}X>{\raggedleft\arraybackslash}X>{\raggedleft\arraybackslash}X>{\raggedleft\arraybackslash}X>{\raggedleft\arraybackslash}X>{\raggedleft\arraybackslash}X}
\toprule
Scale & Dense $\mathrm{Hit@10}$ & Policy $\mathrm{Hit@10}$ mean & Std & 95\% CI & Mean hit delta \\
\midrule
100k & 0.8674 & 0.8652 & 0.0035 & [0.8565, 0.8739] & -0.0022 \\
200k & 0.7970 & 0.8249 & 0.0079 & [0.8052, 0.8446] & +0.0280 \\
400k & 0.7718 & 0.7819 & 0.0044 & [0.7709, 0.7929] & +0.0101 \\
638k & 0.7282 & 0.7466 & 0.0089 & [0.7246, 0.7687] & +0.0185 \\
\bottomrule
\end{tabularx}
\end{table}

\begin{table}[pos=tbp]
\centering
\small
\caption{Multi-seed final context-token stability}
\label{tab:s2}
\begin{tabularx}{\textwidth}{>{\raggedright\arraybackslash}X>{\raggedleft\arraybackslash}X>{\raggedleft\arraybackslash}X>{\raggedleft\arraybackslash}X>{\raggedleft\arraybackslash}X>{\raggedleft\arraybackslash}X>{\raggedleft\arraybackslash}X}
\toprule
Scale & Dense $\mathrm{Tokens@10}$ & Policy $\mathrm{Tokens@10}$ mean & Std & 95\% CI & Mean token saving & Saving 95\% CI \\
\midrule
100k & 1472.39 & 1401.24 & 11.49 & [1372.70, 1429.79] & 4.83\% & [2.89\%, 6.77\%] \\
200k & 1444.12 & 1376.46 & 4.61 & [1365.01, 1387.91] & 4.69\% & [3.89\%, 5.48\%] \\
400k & 1482.30 & 1403.43 & 31.10 & [1326.16, 1480.69] & 5.32\% & [0.11\%, 10.53\%] \\
638k & 1525.62 & 1451.49 & 3.83 & [1441.97, 1461.00] & 4.86\% & [4.24\%, 5.48\%] \\
\bottomrule
\end{tabularx}
\end{table}

The token-saving direction is consistent across scales. The wider 400k token interval should be interpreted as operating-point variance across the routed ranking and context policy, not hidden as a uniform result.

\section{Static Retrieval Baselines}
\label{sec:12-appendix:static-retrieval-baselines}

Dense retrieval is the primary quality baseline. BM25 supplies lexical coverage, but it is weaker as a standalone retriever on LoTTE. Static dense-plus-BM25 reciprocal-rank fusion is competitive at some scales but does not consistently dominate dense. Table~\ref{tab:s3} reports these static baselines.

\begin{table}[pos=tbp]
\centering
\small
\caption{Static LoTTE retrieval baselines across corpus scale}
\label{tab:s3}
\begin{tabularx}{\textwidth}{>{\raggedright\arraybackslash}X>{\raggedleft\arraybackslash}X>{\raggedleft\arraybackslash}X>{\raggedleft\arraybackslash}X>{\raggedleft\arraybackslash}X}
\toprule
Scale & Corpus chunks & BM25 $\mathrm{Hit@10}$ & Dense $\mathrm{Hit@10}$ & Static hybrid $\mathrm{Hit@10}$ \\
\midrule
100k & 101311 & 0.7232 & 0.8674 & 0.8624 \\
200k & 201010 & 0.6292 & 0.7970 & 0.8003 \\
400k & 400674 & 0.5721 & 0.7718 & 0.7617 \\
638k & 638509 & 0.5084 & 0.7282 & 0.7181 \\
\bottomrule
\end{tabularx}
\end{table}

The declining dense score as corpus scale grows motivates adaptive context control, but it does not make dense retrieval obsolete. Dense remains an important recall floor and fallback route in IntentRoute.

\section{Cost Metric Separation}
\label{sec:12-appendix:cost-metric-separation}

The experiments separate three efficiency layers:

\begin{enumerate}
\item source candidate cost: candidates considered before final fusion;
\item dense invocation rate: fraction of queries using global dense retrieval;
\item final context tokens: retrieved chunk tokens sent to the generator.
\end{enumerate}

The main paper claim uses the third layer. Historical routing experiments showed that reducing candidate counts does not automatically reduce final context tokens when the final context remains fixed at top-10. Table~\ref{tab:s4} records the correction audit that motivated this separation.

In Supplementary Table S4, every row labeled \nolinkurl{LoTTE} refers specifically to the technology/search domain; science/search is reported separately in Supplementary Section S8.

\begin{table}[pos=tbp]
\centering
\small
\caption{Representative fixed-top-10 correction audit}
\label{tab:s4}
\begin{tabularx}{\textwidth}{>{\raggedright\arraybackslash}X>{\raggedright\arraybackslash}X>{\raggedleft\arraybackslash}X>{\raggedleft\arraybackslash}X>{\raggedleft\arraybackslash}X>{\raggedleft\arraybackslash}X}
\toprule
Dataset / scale & Routing setting & $\mathrm{Hit@10}$ & Avg tokens@10 & Ratio vs dense & Source candidate cost \\
\midrule
Banking77 & Gated cost-aware routing & 0.9813 & 120.82 & 0.9978x & 142.51 \\
eManual & Gated cost-aware routing & 0.0116 & 17.92 & 0.9829x & 214.07 \\
LoTTE 100k & Quality-first routing & 0.8770 & 1518.44 & 1.0313x & 229.97 \\
LoTTE 100k & Conditional fallback routing & 0.8747 & 1516.24 & 1.0298x & 227.29 \\
LoTTE 100k & Cluster-credit routing & 0.8764 & 1550.65 & 1.0532x & 181.47 \\
LoTTE 200k & Initial gated routing & 0.8154 & 1549.39 & 1.0729x & 232.01 \\
LoTTE 400k & Initial gated routing & 0.7836 & 1547.66 & 1.0441x & 233.22 \\
LoTTE 638k & Initial gated routing & 0.7343 & 1599.95 & 1.0487x & 236.22 \\
\bottomrule
\end{tabularx}
\end{table}

This audit motivated explicit final-context control. The conservative historical policy reduces context size in high-confidence cases, while the stronger main result uses an independently calibrated length budget. Neither candidate-count savings nor route confidence alone establishes prompt-token savings.

\section{Secondary Datasets and Boundary Cases}
\label{sec:12-appendix:secondary-datasets-and-boundary-cases}

The secondary datasets have different roles and should not be pooled into the main LoTTE evidence claim.

\begin{itemize}
\item \textbf{PubMedQA} is an evidence-retrieval proof-of-concept near a dense ceiling. Dense reaches $\mathrm{Hit@10}=0.9930$; trust-weighted feedback reaches $\mathrm{Hit@10}=0.9940$, last reward $0.8727$, and selected-cluster hit $0.8860$. The ground truth is abstract-level context, not a strict answer sentence.
\item \textbf{CovidQA-RAG} is the biomedical transfer row with a non-ceiling dense baseline. The native-full RAGBench checkpoint contains 32,392 chunks and 1,726 evaluated queries after skipping 39 without usable ground truth. Dense reaches $\mathrm{Hit@10}=0.6095$, BM25 reaches $0.4884$, hybrid RRF reaches $0.6037$, no-feedback IntentRoute reaches $0.6294$, and trust-weighted IntentRoute reaches $0.6300$ at fixed top-10. In the five-fold budgeted protocol, four of five folds select a compressed route, mean hit delta is -0.21 percentage points versus dense, mean final-context token saving is 8.34\%, and strict 1pp non-inferiority remains $0/3$ seeds.
\item \textbf{Banking77} is an intent-routing proxy rather than an evidence-retrieval benchmark. Dense/reference $\mathrm{Hit@10}$ is $0.9805$; trust-weighted feedback reaches $\mathrm{Hit@10}=0.9844$, last reward $0.9805$, and selected-cluster hit $0.9983$. It supports the feedback mechanism, not the main evidence-retrieval headline.
\item \textbf{eManual} is a duplicate-text limitation case. Strict dense $\mathrm{Hit@10}$ is $0.3231$, text-equivalent dense $\mathrm{Hit@10}$ is $0.5615$, and deduplicated dense $\mathrm{Hit@10}$ rises to $0.8615$. Strict chunk IDs can therefore understate useful retrieval.
\item \textbf{CUAD} is a sparse legal smoke/stress case. Dense/reference $\mathrm{Hit@10}$ is $0.0759$ and the trust-weighted smoke reaches $\mathrm{Hit@10}=0.0886$. It is a GT-anchored sample, not positive full-corpus evidence.
\end{itemize}

\subsection{eManual Duplicate-Text Diagnostic}
\label{sec:12-appendix:emanual-duplicate-text-diagnostic}

eManual contains 18,812 corpus chunks but only 1,729 unique text strings. Strict chunk-id evaluation can therefore mark semantically equivalent retrievals as incorrect. Table~\ref{tab:s5} quantifies the strict, text-equivalent, and deduplicated views.

\begin{table}[pos=tbp]
\centering
\small
\caption{eManual strict, text-equivalent, and deduplicated evaluation}
\label{tab:s5}
\begin{tabularx}{\textwidth}{>{\raggedright\arraybackslash}X>{\raggedright\arraybackslash}X>{\raggedleft\arraybackslash}X>{\raggedleft\arraybackslash}X>{\raggedleft\arraybackslash}X}
\toprule
Method & Evaluation mode & $\mathrm{Hit@10}$ & $\mathrm{MRR@10}$ & $\mathrm{nDCG@10}$ \\
\midrule
BM25 & Strict chunk ID & 0.1154 & 0.0244 & 0.0256 \\
BM25 & Text-equivalent & 0.3846 & 0.3059 & 0.1620 \\
Dense & Strict chunk ID & 0.3231 & 0.0551 & 0.0526 \\
Dense & Text-equivalent & 0.5615 & 0.4716 & 0.2030 \\
Static hybrid & Strict chunk ID & 0.1692 & 0.0366 & 0.0287 \\
Static hybrid & Text-equivalent & 0.5846 & 0.4895 & 0.2263 \\
Dense & Deduplicated corpus & 0.8615 & 0.5736 & 0.3807 \\
\bottomrule
\end{tabularx}
\end{table}

The text-equivalent and deduplicated metrics are diagnostics, not replacements for the strict evaluation. They show that eManual's low strict score cannot be interpreted as proof that useful local structure is absent.

\section{Encoder Robustness}
\label{sec:12-appendix:encoder-robustness}

The main scale-up uses \nolinkurl{sentence-transformers/all-MiniLM-L6-v2}. A LoTTE technology/search 100k robustness check replaces it with \nolinkurl{sentence-transformers/multi-qa-MiniLM-L6-cos-v1}, a QA-tuned MiniLM-family encoder with the same 384-dimensional embedding size and a similar CPU-friendly resource class. Tables~\ref{tab:s6} and~\ref{tab:s7} report encoder-family and matched-backbone robustness.

\begin{table}[pos=tbp]
\centering
\small
\caption{QA-tuned MiniLM-family encoder robustness}
\label{tab:s6}
\begin{tabularx}{\textwidth}{>{\raggedright\arraybackslash}X>{\raggedleft\arraybackslash}X>{\raggedleft\arraybackslash}X>{\raggedleft\arraybackslash}X>{\raggedleft\arraybackslash}X>{\raggedleft\arraybackslash}X>{\raggedleft\arraybackslash}X}
\toprule
Method & $\mathrm{Hit@10}$ & $\mathrm{MRR@10}$ & $\mathrm{nDCG@10}$ & Evid. recall@10 & Avg tokens@10 & Token ratio \\
\midrule
Dense-only & 0.8809 & 0.7220 & 0.6616 & 0.7163 & 1514.51 & 1.0000x \\
Conservative policy & 0.8853 & 0.7118 & 0.6291 & 0.6789 & 1463.71 & 0.9665x \\
\bottomrule
\end{tabularx}
\end{table}

Under the QA-tuned encoder, the dense baseline becomes stronger and the conservative policy still preserves dense-level query hit while reducing final context tokens by 3.35\%. Ranking metrics and evidence recall are lower than dense, so this is a bounded robustness result rather than a universal retrieval-metric improvement.

\begin{table}[pos=tbp]
\centering
\small
\caption{Matched-backbone context-budget robustness on the frozen LoTTE technology/search 100k split}
\label{tab:s7}
\begin{tabularx}{\textwidth}{>{\raggedright\arraybackslash}X>{\raggedright\arraybackslash}X>{\raggedleft\arraybackslash}X>{\raggedleft\arraybackslash}X>{\raggedright\arraybackslash}X>{\raggedleft\arraybackslash}X}
\toprule
Backbone & Route mode & Dense $\mathrm{Hit@10}$ & Method $\mathrm{Hit@10}$ & Hit delta & Token saving \\
\midrule
MiniLM & calibrated multi-route & 0.8705 & 0.8705 & +0.00 pp & 6.18\% \\
BGE-base & full multi-route & 0.8993 & 0.8985 & -0.08 pp & 11.99\% \\
E5-base & full multi-route & 0.8753 & 0.8689 & -0.64 pp & 12.20\% \\
BGE-base & quality-first & 0.8993 & 0.9081 & +0.88 pp & 7.23\% \\
\bottomrule
\end{tabularx}
\end{table}

The full multi-route rows show that the quality-cost pattern is not tied to the MiniLM backbone. More aggressive gated BGE/E5 variants lose more hit and are treated as boundary settings. The BGE quality-first row demonstrates that the frontier is tunable; an equivalent above-dense E5 point was not found on this split.

\section{Downstream Answer-Level Evaluation}
\label{sec:12-appendix:downstream-answer-level-evaluation}

The formal downstream evaluation uses 300 deterministic queries from the frozen LoTTE technology/search 100k test split. Seven methods produce 2,100 answers with \nolinkurl{deepseek-v4-flash}. The fixed answers receive 2,100 DeepSeek, 2,100 GLM-5.2, and 2,065 MiniMax-M3 schema-valid judgments. Cross-judge statistics use the 2,065 query-method keys shared by all judges. Tables~\ref{tab:s8}, \ref{tab:s9}, \ref{tab:s10}, \ref{tab:s11}, and~\ref{tab:s12} report method results, paired tests, judge coverage, agreement, and majority-vote comparisons.

\begin{table}[pos=tbp]
\centering
\small
\caption{DeepSeek-judged downstream answer and context results}
\label{tab:s8}
\begin{tabularx}{\textwidth}{>{\raggedright\arraybackslash}X>{\raggedleft\arraybackslash}X>{\raggedleft\arraybackslash}X>{\raggedleft\arraybackslash}X>{\raggedleft\arraybackslash}X>{\raggedleft\arraybackslash}X>{\raggedleft\arraybackslash}X}
\toprule
Method & Correct & Faithful & Strict citation support & Insufficient context & Avg context tokens & Tokens / correct \\
\midrule
MiniLM dense & 0.9200 & 0.9533 & 0.3533 & 0.0967 & 1461 & 1588 \\
BGE dense & 0.9167 & 0.9433 & 0.3567 & 0.0767 & 1698 & 1852 \\
BGE IntentRoute & 0.9167 & 0.9200 & 0.3700 & 0.0767 & 1596 & 1741 \\
E5 dense & 0.9167 & 0.9300 & 0.4100 & 0.0700 & 1525 & 1663 \\
E5 IntentRoute & 0.9200 & 0.9333 & 0.3633 & 0.0567 & 1341 & 1458 \\
Dense+MMR & 0.8900 & 0.9100 & 0.0733 & 0.0800 & 1240 & 1393 \\
IntentRoute + MMR & 0.9133 & 0.9267 & 0.0833 & 0.0900 & 1157 & 1267 \\
\bottomrule
\end{tabularx}
\end{table}

\begin{table}[pos=tbp]
\centering
\small
\caption{Original DeepSeek-judged paired downstream comparisons}
\label{tab:s9}
\begin{tabularx}{\textwidth}{>{\raggedright\arraybackslash}X>{\raggedright\arraybackslash}X>{\raggedright\arraybackslash}X>{\raggedleft\arraybackslash}X>{\raggedleft\arraybackslash}X>{\raggedleft\arraybackslash}X}
\toprule
Comparison & Correct delta & 95\% CI & McNemar $p$ & Token saving & 95\% CI \\
\midrule
BGE IntentRoute vs dense & +0.00 pp & [-2.67, +2.67] pp & 1.000 & 6.00\% & [4.01\%, 7.97\%] \\
E5 IntentRoute vs dense & +0.33 pp & [-3.00, +3.67] pp & 1.000 & 12.04\% & [9.93\%, 14.16\%] \\
IntentRoute + MMR vs Dense + MMR & +2.33 pp & [-1.67, +6.33] pp & 0.324 & 6.65\% & [4.28\%, 8.97\%] \\
\bottomrule
\end{tabularx}
\end{table}

The context-saving intervals are positive while correctness intervals include zero. The multi-judge extension below tests whether this conclusion depends on the original DeepSeek judge.

\begin{table}[pos=tbp]
\centering
\small
\caption{Multi-judge coverage and calibration}
\label{tab:s10}
\begin{tabularx}{\textwidth}{>{\raggedright\arraybackslash}X>{\raggedleft\arraybackslash}X>{\raggedleft\arraybackslash}X>{\raggedleft\arraybackslash}X>{\raggedleft\arraybackslash}X>{\raggedleft\arraybackslash}X>{\raggedleft\arraybackslash}X>{\raggedleft\arraybackslash}X}
\toprule
Judge & Valid & Coverage & Correctness mean & Correct & Faithfulness mean & Faithful & Citations supported \\
\midrule
DeepSeek & 2,100 & 100.00\% & 4.672 & 91.33\% & 4.782 & 93.10\% & 89.57\% \\
GLM-5.2 & 2,100 & 100.00\% & 4.550 & 88.24\% & 4.734 & 92.81\% & 92.14\% \\
MiniMax-M3 & 2,065 & 98.33\% & 4.293 & 85.71\% & 4.633 & 95.45\% & 94.48\% \\
\bottomrule
\end{tabularx}
\end{table}

MiniMax-M3 rejects 35 query-method inputs spanning 18 queries through provider-side content filtering. These values are not imputed. Absolute score calibration differs across judges, so raw ordinal scores are not pooled.

\begin{table}[pos=tbp]
\centering
\small
\caption{Pairwise agreement on 2,065 shared judgments}
\label{tab:s11}
\begin{tabular}{llrr}
\toprule
Field & Judge pair & Raw agreement & Cohen's $\kappa$ \\
\midrule
Correct & DS / GLM & 91.04\% & 0.508 \\
Correct & DS / MM & 89.88\% & 0.503 \\
Correct & GLM / MM & 92.15\% & 0.653 \\
Faithful & DS / GLM & 91.72\% & 0.364 \\
Faithful & DS / MM & 93.27\% & 0.374 \\
Faithful & GLM / MM & 93.41\% & 0.405 \\
\bottomrule
\end{tabular}
\end{table}

Here DS denotes DeepSeek and MM denotes MiniMax-M3. Three-judge unanimity is 86.54\% for correctness and 89.20\% for faithfulness. The corresponding majority-positive rates are 88.96\% and 95.35\%. High raw faithfulness agreement coexists with lower $\kappa$ because positive judgments are highly prevalent.

\begin{table}[pos=tbp]
\centering
\small
\caption{Three-judge-majority paired comparisons}
\label{tab:s12}
\begin{tabularx}{\textwidth}{>{\raggedright\arraybackslash}X>{\raggedleft\arraybackslash}X>{\raggedright\arraybackslash}X>{\raggedleft\arraybackslash}X>{\raggedright\arraybackslash}X>{\raggedleft\arraybackslash}X>{\raggedleft\arraybackslash}X}
\toprule
Comparison & $n$ & Correct delta (95\% CI) & McNemar $p$ & Faithful delta (95\% CI) & McNemar $p$ & Context saving (95\% CI) \\
\midrule
BGE IntentRoute vs dense & 289 & -3.46 pp [-6.92, 0.00] & 0.0755 & -4.15 pp [-6.92, -1.73] & 0.0018 & 6.27\% [4.22\%, 8.26\%] \\
E5 IntentRoute vs dense & 289 & -2.08 pp [-5.88, +1.73] & 0.3616 & -0.69 pp [-3.81, +2.42] & 0.8238 & 11.97\% [9.78\%, 14.17\%] \\
IntentRoute + MMR vs Dense + MMR & 295 & +0.34 pp [-3.39, +4.07] & 1.0000 & +4.07 pp [+0.68, +7.46] & 0.0290 & 6.75\% [4.40\%, 9.11\%] \\
\bottomrule
\end{tabularx}
\end{table}

No individual judge or majority-vote comparison finds a significant correctness difference. This supports a bounded correctness-robustness claim, not strict non-inferiority. Faithfulness is mixed: the BGE majority result is negative, while the SentMMR composition is positive. The evaluation remains LLM-as-judge evidence rather than human evaluation. The under-specified \nolinkurl{insufficient_context_appropriate} field is retained in raw artifacts but excluded from headline analysis.

\section{Calibration/Test Context-Budget Validation}
\label{sec:12-appendix:calibration-test-context-budget-validation}

The calibration/test protocol selects the final-context budget on calibration queries and freezes it before evaluation on held-out test queries. Tables~\ref{tab:s13}, \ref{tab:s14}, \ref{tab:s15}, and~\ref{tab:s16} report the frozen split, independent calibration, partition sensitivity, and normalized five-fold audit.

\begin{table}[pos=tbp]
\centering
\small
\caption{Frozen context-budget validation on LoTTE technology/search}
\label{tab:s13}
\begin{tabularx}{\textwidth}{>{\raggedright\arraybackslash}X>{\raggedright\arraybackslash}X>{\raggedright\arraybackslash}X>{\raggedright\arraybackslash}X>{\raggedleft\arraybackslash}X>{\raggedright\arraybackslash}X>{\raggedleft\arraybackslash}X}
\toprule
Scale & Selected policy & Calibration eligible & Hit delta vs dense & Token saving & Dense trunc. hit delta & Dense trunc. token saving \\
\midrule
100k & \nolinkurl{token_budget_r0.95_m4} & True & +0.00 pp & 6.18\% & -1.44 pp & 13.83\% \\
200k & \nolinkurl{token_budget_r0.85_m4} & True & +1.20 pp & 16.00\% & -2.40 pp & 21.95\% \\
400k & \nolinkurl{token_budget_r0.98_m4} & False / original split & +2.32 pp & 6.57\% & -0.24 pp & 11.44\% \\
638k & \nolinkurl{token_budget_r0.85_m4} & True & -0.08 pp & 17.53\% & -3.84 pp & 21.90\% \\
\bottomrule
\end{tabularx}
\end{table}

The calibrated policies should be compared against dense-only adaptive truncation because both reduce final context size. IntentRoute preserves substantially more $\mathrm{Hit@10}$ at a still meaningful token saving level. The 400k row is diagnostic rather than calibration-eligible in the original artifact set. Supplementary Table S16 reports the completed cross-fitted follow-up.

\begin{table}[pos=tbp]
\centering
\small
\caption{Independently calibrated 100k quality constraints}
\label{tab:s14}
\begin{tabularx}{\textwidth}{>{\raggedright\arraybackslash}X>{\raggedright\arraybackslash}X>{\raggedright\arraybackslash}X>{\raggedleft\arraybackslash}X>{\raggedright\arraybackslash}X>{\raggedright\arraybackslash}X>{\raggedleft\arraybackslash}X}
\toprule
Calibration Hit margin & IntentRoute policy & IR test Hit delta & IR saving & Dense policy & Dense test Hit delta & Dense saving \\
\midrule
0.0pp & \nolinkurl{r0.95/m4} & +0.00pp & 6.18\% & \nolinkurl{r1.00/m4} & +0.00pp & 0.00\% \\
0.5pp & \nolinkurl{r0.93/m4} & -0.32pp & 7.97\% & \nolinkurl{r1.00/m4} & +0.00pp & 0.00\% \\
1.0pp & \nolinkurl{r0.84/m4} & -1.28pp & 17.13\% & \nolinkurl{r1.00/m4} & +0.00pp & 0.00\% \\
2.0pp & \nolinkurl{r0.80/m4} & -1.60pp & 20.90\% & \nolinkurl{r0.82/m4} & -2.40pp & 25.22\% \\
\bottomrule
\end{tabularx}
\end{table}

The zero-margin row preserves equal mean test Hit while selecting nonzero IntentRoute saving, but strict seed-level non-inferiority remains 0/3. Held-out same-saving interpolation is descriptive only and finds small IntentRoute-minus- Dense Hit differences from \nolinkurl{+0.47pp} at 5\% saving to \nolinkurl{-0.01pp} at 20\%.

\begin{table}[pos=tbp]
\centering
\small
\caption{Calibration-partition sensitivity over 20 overlapping splits}
\label{tab:s15}
\begin{tabularx}{\textwidth}{>{\raggedright\arraybackslash}X>{\raggedright\arraybackslash}X>{\raggedright\arraybackslash}X>{\raggedright\arraybackslash}X>{\raggedleft\arraybackslash}X>{\raggedright\arraybackslash}X}
\toprule
Scale & Eligible splits & Test Hit range & Mean test Hit delta & Saving range & Within 1pp of dense \\
\midrule
100k & 12/20 & [-2.00, +0.72]pp & -0.45pp & [6.18\%, 17.73\%] & 14/20 \\
200k & 19/20 & [+0.56, +3.52]pp & +1.53pp & [5.33\%, 17.29\%] & 20/20 \\
400k & 16/20 & [-2.08, +2.80]pp & +0.45pp & [6.57\%, 18.27\%] & 17/20 \\
638k & 19/20 & [-0.88, +2.24]pp & +0.44pp & [7.85\%, 17.91\%] & 20/20 \\
\bottomrule
\end{tabularx}
\end{table}

These partitions reuse the same frozen rankings and overlap in their query membership. They diagnose policy-selection sensitivity and must not be counted as 20 independent experiments. The result supports stronger split stability at 200k/638k, moderate sensitivity at 100k, and continued diagnostic treatment of 400k.

\begin{table}[pos=tbp]
\centering
\small
\caption{Normalized five-fold out-of-fold calibration using identical canonical query folds and policy rules across scales}
\label{tab:s16}
\begin{tabularx}{\textwidth}{>{\raggedright\arraybackslash}X>{\raggedright\arraybackslash}X>{\raggedright\arraybackslash}X>{\raggedleft\arraybackslash}X>{\raggedright\arraybackslash}X>{\raggedleft\arraybackslash}X>{\raggedright\arraybackslash}X}
\toprule
Scale & Eligible folds & Mean Hit delta & Mean token saving & Strict NI seeds & Selected-policy count & Dense compressed folds \\
\midrule
100k & 2/5 & -1.06pp & 4.16\% & 0/3 & 2 & 0/5 \\
200k & 5/5 & +1.40pp & 16.07\% & 2/3 & 1 & 0/5 \\
400k & 5/5 & +0.00pp & 14.50\% & 0/3 & 5 & 0/5 \\
638k & 5/5 & +0.28pp & 15.23\% & 0/3 & 3 & 0/5 \\
\bottomrule
\end{tabularx}
\end{table}

Each canonical LoTTE query is held out exactly once and remains in the same fold at every corpus scale. Dense and IntentRoute independently select from the predefined budget grid; when no policy satisfies the zero-drop gate, the method uses Dense top-10 fallback. The 400k result closes the missing normalized calibration check, but its five distinct policies and 0/3 strict non-inferiority result retain the partition-sensitivity boundary. The 100k row similarly shows that positive original-split behavior does not imply uniform cross-fitted behavior.

\section{Cross-Domain Validation}
\label{sec:12-appendix:cross-domain-validation}

LoTTE science/search is used as a second-domain validation, not as a replacement for the main LoTTE technology/search scale-up. Tables~\ref{tab:s17} and~\ref{tab:s18} separate ranking transfer from frozen-budget behavior.

\begin{table}[pos=tbp]
\centering
\small
\caption{Science/search fixed top-10 ranking validation}
\label{tab:s17}
\begin{tabularx}{\textwidth}{>{\raggedright\arraybackslash}X>{\raggedleft\arraybackslash}X>{\raggedleft\arraybackslash}X>{\raggedleft\arraybackslash}X>{\raggedleft\arraybackslash}X>{\raggedright\arraybackslash}X}
\toprule
Domain/scale & Corpus chunks & Queries & Dense $\mathrm{Hit@10}$ & IntentRoute $\mathrm{Hit@10}$ & Hit delta \\
\midrule
science/search 20k/q200 & 20,490 & 200 & 0.8950 & 0.9267 & +3.17 pp \\
science/search 100k & 101,187 & 596 & 0.8926 & 0.9077 & +1.51 pp \\
\bottomrule
\end{tabularx}
\end{table}

\begin{table}[pos=tbp]
\centering
\small
\caption{Science/search frozen context-budget validation}
\label{tab:s18}
\begin{tabularx}{\textwidth}{>{\raggedright\arraybackslash}X>{\raggedright\arraybackslash}X>{\raggedleft\arraybackslash}X>{\raggedright\arraybackslash}X>{\raggedleft\arraybackslash}X>{\raggedright\arraybackslash}X}
\toprule
Domain/scale & Budget policy & Seed & Frozen test hit delta vs dense & Token saving & Strict NI by CI \\
\midrule
20k/q200 & \nolinkurl{token_budget_r0.85_m4} & 13 & +2.86 pp & 13.18\% & True \\
20k/q200 & \nolinkurl{token_budget_r0.85_m4} & 17 & +1.43 pp & 14.31\% & False \\
20k/q200 & \nolinkurl{token_budget_r0.85_m4} & 19 & +0.71 pp & 13.91\% & False \\
100k & \nolinkurl{token_budget_r0.85_m4} & 13 & -1.20 pp & 19.21\% & False \\
100k & \nolinkurl{token_budget_r0.85_m4} & 17 & +0.00 pp & 17.53\% & False \\
100k & \nolinkurl{token_budget_r0.85_m4} & 19 & -0.96 pp & 20.53\% & False \\
\bottomrule
\end{tabularx}
\end{table}

The fixed top-10 ranking gains transfer, while aggressive final-context budgets require domain and scale calibration.

\section{Feedback-Driven Hard-Case Recovery}
\label{sec:12-appendix:feedback-driven-hard-case-recovery}

Hard-case recovery focuses on affected queries where dense top-10 retrieves at least one GT chunk but the budgeted IntentRoute context misses. Tables~\ref{tab:s19} and~\ref{tab:s20} distinguish same-query repair from held-out recovery.

\begin{table}[pos=tbp]
\centering
\small
\caption{Same-query feedback recovery on affected queries}
\label{tab:s19}
\begin{tabularx}{\textwidth}{>{\raggedright\arraybackslash}X>{\raggedright\arraybackslash}X>{\raggedleft\arraybackslash}X>{\raggedleft\arraybackslash}X>{\raggedleft\arraybackslash}X>{\raggedleft\arraybackslash}X}
\toprule
Domain & Retry method & Affected queries & Recovered & Recovery rate & Avg token saving \\
\midrule
science 100k & arm boost & 34 & 5 & 14.71\% & 17.40\% \\
science 100k & arm boost + conservative budget & 34 & 14 & 41.18\% & 5.76\% \\
science 100k & full-context fallback & 34 & 17 & 50.00\% & -8.07\% \\
technology 100k & arm boost & 42 & 8 & 19.05\% & 13.68\% \\
technology 100k & arm boost + conservative budget & 42 & 9 & 21.43\% & 11.75\% \\
technology 100k & full-context fallback & 42 & 12 & 28.57\% & 0.96\% \\
\bottomrule
\end{tabularx}
\end{table}

Same-query retry is post-feedback repair evidence. It is not a first-pass generalization result.

\begin{table}[pos=tbp]
\centering
\small
\caption{Calibration-to-test recovery generalization}
\label{tab:s20}
\begin{tabularx}{\textwidth}{>{\raggedright\arraybackslash}X>{\raggedright\arraybackslash}X>{\raggedright\arraybackslash}X>{\raggedleft\arraybackslash}X}
\toprule
Domain & Recovery policy & Mean hit delta & Avg token saving \\
\midrule
science 100k & conservative budget on learned risky arms after calibration & +0.16 pp & 16.13\% \\
science 100k & full-context fallback on learned risky arms after calibration & +0.48 pp & 13.09\% \\
technology 100k & conservative budget on learned risky arms after calibration & -0.16 pp & 5.88\% \\
technology 100k & full-context fallback on learned risky arms after calibration & +0.16 pp & 4.25\% \\
\bottomrule
\end{tabularx}
\end{table}

The held-out effect is small and domain-dependent. Feedback should therefore be used as a controlled fallback trigger rather than as unconditional global reranking.

\subsection{Frozen Unseen-Query Policy Audit}
\label{sec:12-appendix:frozen-unseen-query-policy-audit}

The formal frozen-policy audit trains learned route state on four disjoint canonical folds for eight prequential epochs, freezes policy matrices, feedback memory, route statistics, and reward history, then ranks the fifth fold once with zero held-out feedback updates. This retrieval-only audit covers 596 queries per domain, five folds, and seeds 13/17/19.

On technology/search 100k, learned full routing reaches $\mathrm{Hit@10}=0.8792$ versus 0.8674 for Dense, but is 0.11pp below cold no-feedback full routing and has no significant paired advantage over static-nearest full routing. On science/search 100k, learned full routing reaches 0.9004 versus 0.8926 for Dense, but is 0.39pp below cold no-feedback full and 0.56pp below static-nearest full. Learned gated routing is significantly below Dense in all seeds in both domains (-4.08pp and -5.59pp mean deltas). The audit therefore supports transfer of the full Dense/BM25-rescued route surface, not a universal first-pass advantage of learned simulated feedback.

\section{Strong Post-Retrieval Baselines}
\label{sec:12-appendix:strong-post-retrieval-baselines}

These baselines test whether simpler post-retrieval operations explain the main final-context result. Tables~\ref{tab:s21}, \ref{tab:s22}, \ref{tab:s23}, and~\ref{tab:s24} report matched compression, reranking, and prompt-pruning controls.

\begin{table}[pos=tbp]
\centering
\small
\caption{Dense+Sentence-MMR same-budget baseline on LoTTE technology/search 100k}
\label{tab:s21}
\begin{tabularx}{\textwidth}{>{\raggedright\arraybackslash}X>{\raggedleft\arraybackslash}X>{\raggedleft\arraybackslash}X>{\raggedleft\arraybackslash}X>{\raggedleft\arraybackslash}X>{\raggedleft\arraybackslash}X}
\toprule
Budget target & $\mathrm{Hit@10}$ & Hit delta vs dense & Evid. recall@10 & Avg context tokens & Token saving vs dense \\
\midrule
Dense top-10 & 0.8705 & 0.0000 & 0.7081 & 1470.1 & 0.00\% \\
SentMMR seed13 budget & 0.8705 & 0.0000 & 0.7081 & 1287.8 & 12.40\% \\
SentMMR seed17 budget & 0.8705 & 0.0000 & 0.7075 & 1278.0 & 13.07\% \\
SentMMR seed19 budget & 0.8705 & 0.0000 & 0.7081 & 1302.1 & 11.43\% \\
\bottomrule
\end{tabularx}
\end{table}

Dense+Sentence-MMR preserves dense chunk-support on this split while reducing selected sentence tokens. It is therefore a strong final-context compression baseline, not a weak control.

\begin{table}[pos=tbp]
\centering
\small
\caption{Compressor-normalized comparison on LoTTE technology/search 100k}
\label{tab:s22}
\begin{tabularx}{\textwidth}{>{\raggedright\arraybackslash}X>{\raggedright\arraybackslash}X>{\raggedleft\arraybackslash}X>{\raggedleft\arraybackslash}X>{\raggedleft\arraybackslash}X}
\toprule
Source pool & Method & Ratio & $\mathrm{Hit@10}$ range & Saving vs dense \\
\midrule
Dense & top-10 & - & 0.8705 & 0.00\% \\
Dense & SentMMR & 0.95 & 0.8705 & 5.33\% \\
Dense & SentMMR & 0.90 & 0.8705 & 10.22\% \\
Dense & SentMMR & 0.85 & 0.8705 & 15.16\% \\
IntentRoute & target & - & 0.8657-0.8777 & 4.98-7.14\% \\
IntentRoute & SentMMR & 0.95 & 0.8657-0.8777 & 10.07-12.14\% \\
IntentRoute & SentMMR & 0.90 & 0.8657-0.8777 & 14.72-16.68\% \\
IntentRoute & SentMMR & 0.85 & 0.8657-0.8777 & 19.41-21.24\% \\
\bottomrule
\end{tabularx}
\end{table}

Applying the same compressor to dense and IntentRoute evidence pools supports the route-and-budget controller framing. The compressor is shared; the evidence pool and budget controller determine the starting point.

\begin{table}[pos=tbp]
\centering
\small
\caption{Cross-encoder reranker baseline on LoTTE technology/search 100k}
\label{tab:s23}
\begin{tabularx}{\textwidth}{>{\raggedright\arraybackslash}X>{\raggedleft\arraybackslash}X>{\raggedleft\arraybackslash}X>{\raggedleft\arraybackslash}X>{\raggedleft\arraybackslash}X}
\toprule
Method & $\mathrm{Hit@10}$ & Evid. recall@10 & Avg context tokens & Token saving vs dense \\
\midrule
Dense top-10 & 0.8705 & 0.7081 & 1470 & 0.00\% \\
Cross-encoder top-10 & 0.8777 & 0.7332 & 1792 & -21.91\% \\
IntentRoute target & 0.8657-0.8777 & 0.6766-0.6871 & 1365-1397 & 4.98-7.14\% \\
Cross-encoder same budget & 0.8633-0.8729 & 0.6975-0.7044 & 1360-1390 & 5.43-7.49\% \\
\bottomrule
\end{tabularx}
\end{table}

The cross-encoder reranker improves full top-10 support metrics, but that full reranked context is longer on average. Under the same per-query token budgets as IntentRoute, reranking does not uniformly dominate the calibrated controller.

\begin{table}[pos=tbp]
\centering
\small
\caption{SelectiveContext-lite prompt-pruning baseline}
\label{tab:s24}
\begin{tabularx}{\textwidth}{>{\raggedright\arraybackslash}X>{\raggedleft\arraybackslash}X>{\raggedleft\arraybackslash}X>{\raggedleft\arraybackslash}X>{\raggedleft\arraybackslash}X}
\toprule
Source pool & Ratio & $\mathrm{Hit@10}$ & Token saving vs dense & Extra saving vs source \\
\midrule
Dense & 0.95 & 0.8705 & 5.66\% & 5.66\% \\
Dense & 0.90 & 0.8705 & 10.42\% & 10.42\% \\
Dense & 0.85 & 0.8705 & 15.31\% & 15.31\% \\
Dense & 0.75 & 0.8705 & 25.19\% & 25.19\% \\
IntentRoute & 0.95 & 0.8657-0.8777 & 10.38-12.42\% & 5.62-5.69\% \\
IntentRoute & 0.90 & 0.8657-0.8777 & 14.92-16.87\% & 10.46-10.48\% \\
IntentRoute & 0.85 & 0.8657-0.8777 & 19.53-21.40\% & 15.31-15.35\% \\
IntentRoute & 0.75 & 0.8657-0.8777 & 28.95-30.57\% & 25.20-25.23\% \\
\bottomrule
\end{tabularx}
\end{table}

SelectiveContext-lite is a deterministic local proxy, not LLMLingua. Its role is to show that prompt pruning can be stacked after either evidence pool and does not replace upstream route control or final-budget calibration.

\section{Route-Control Attribution and Arm Sensitivity}
\label{sec:12-appendix:route-control-attribution-and-arm-sensitivity}

Tables~\ref{tab:s25}, \ref{tab:s26}, \ref{tab:s27}, and~\ref{tab:s28} isolate geometry, feedback, arm granularity, and frozen-trajectory route mediation.

\begin{table}[pos=tbp]
\centering
\small
\caption{Static geometry versus uniform-random route control}
\label{tab:s25}
\begin{tabularx}{\textwidth}{>{\raggedright\arraybackslash}X>{\raggedleft\arraybackslash}X>{\raggedleft\arraybackslash}X>{\raggedleft\arraybackslash}X>{\raggedright\arraybackslash}X>{\raggedleft\arraybackslash}X}
\toprule
Setting & Full top-10 hit & Route reward & Selected-cluster hit & Test hit delta & Token saving \\
\midrule
Static nearest geometry & 0.8764 & 0.8563 & 0.8870 & +1.44 pp & 5.03\% \\
Uniform random control & 0.8842 & 0.1499 & 0.1577 & +1.04 pp & 11.92\% \\
\bottomrule
\end{tabularx}
\end{table}

Dense/BM25 rescue keeps final fused hit high in both rows, while route reward and selected-cluster hit separate meaningful local routing from random arm selection. Geometry is therefore supported as a route-control signal, not a standalone explanation of final fused quality.

\begin{table}[pos=tbp]
\centering
\small
\caption{Feedback and static route controls}
\label{tab:s26}
\begin{tabularx}{\textwidth}{>{\raggedright\arraybackslash}X>{\raggedleft\arraybackslash}X>{\raggedleft\arraybackslash}X>{\raggedleft\arraybackslash}X>{\raggedright\arraybackslash}X>{\raggedleft\arraybackslash}X}
\toprule
Setting & Route reward & Selected-cluster hit & Dense rate & Test hit delta & Token saving \\
\midrule
Learned full multi-route & 0.6790 & 0.5766 & 1.0000 & -1.68 pp & 17.86\% \\
Learned gated & 0.6790 & 0.5766 & 0.7377 & -5.20 pp & 11.83\% \\
Static nearest gated & 0.8563 & 0.8870 & 0.9586 & -2.40 pp & 12.01\% \\
No-feedback gated & 0.1504 & 0.1570 & 1.0000 & -1.60 pp & 16.56\% \\
\bottomrule
\end{tabularx}
\end{table}

Feedback-updated LinUCB improves route quality over no-feedback/random controls, but the learned gated threshold is a cost-aggressive boundary. Static geometry remains a strong prior, while dense fallback explains part of the fused result.

\begin{table}[pos=tbp]
\centering
\small
\caption{Arm-count sensitivity}
\label{tab:s27}
\begin{tabularx}{\textwidth}{>{\raggedleft\arraybackslash}X>{\raggedleft\arraybackslash}X>{\raggedright\arraybackslash}X>{\raggedleft\arraybackslash}X>{\raggedleft\arraybackslash}X>{\raggedright\arraybackslash}X}
\toprule
$K$ & Static route reward & Full test hit delta & Full token saving & Gated dense rate & Gated hit delta \\
\midrule
8 & 0.9128 & +1.44 pp & 6.23\% & 0.4083 & -1.84 pp \\
16 & 0.8826 & +0.80 pp & 10.49\% & 0.6089 & -1.68 pp \\
32 & 0.8563 & +0.56 pp & 4.68\% & 0.7377 & -4.48 pp \\
64 & 0.8272 & +0.40 pp & 11.19\% & 0.8986 & -3.76 pp \\
128 & 0.8479 & +1.20 pp & 10.23\% & 0.9502 & -3.12 pp \\
\bottomrule
\end{tabularx}
\end{table}

Full multi-route quality is stable across the tested grid, whereas gated dense-saving behavior depends on arm granularity. Cross-scale correlations between geometry diagnostics and final quality-cost gain are mixed and small-sample; final behavior belongs to the complete calibrated controller.

\begin{table}[pos=tbp]
\centering
\small
\caption{Frozen-trajectory dynamic route mediation; Save is relative to uncompressed dense}
\label{tab:s28}
\begin{tabular}{lrrr}
\toprule
Route & Src. Hit & Budget Hit & Save \\
\midrule
Gated & 0.8793 & 0.8705 & 6.18\% \\
Full & 0.8841 & 0.8745 & 5.27\% \\
Shuffled & 0.8313 & 0.8225 & 6.54\% \\
Cluster & 0.7698 & 0.7626 & 6.93\% \\
Dense & 0.8705 & 0.8561 & 13.83\% \\
\bottomrule
\end{tabular}
\end{table}

The replay freezes selected arms and feedback state, and exactly reproduces the original dynamic ranking before changing route shapes. Dynamic gating exceeds the shuffled-tier control by 4.80 percentage points before and after the common budget, with 3/3 paired intervals excluding zero. It does not exceed fixed full fusion; instead, it exposes a bounded quality-cost trade-off. Route confidence has no detected association with oracle safe-token headroom, so the result supports route assignment rather than direct compression-safety prediction.

\section{Reproducibility and Reporting Guardrails}
\label{sec:12-appendix:reproducibility-and-reporting-guardrails}

\subsection{Standard Ranking Metrics}
\label{sec:12-appendix:standard-ranking-metrics}

For mean reciprocal rank, let $\rho_t$ be the rank of the first relevant chunk within the top-$K$ list, or $\infty$ if no relevant chunk is retrieved:

\[
\mathrm{MRR@K}(q_t) =
\begin{cases}
\frac{1}{\rho_t}, & \rho_t \le K, \\
0, & \rho_t = \infty.
\end{cases}
\]

For binary relevance $\mathrm{rel}_{t,j} \in \{0,1\}$, we compute:

\[
\mathrm{DCG@K}(q_t) =
\sum_{j=1}^{K} \frac{\mathrm{rel}_{t,j}}{\log_2(j+1)},
\]

\[
\mathrm{nDCG@K}(q_t) =
\frac{\mathrm{DCG@K}(q_t)}{\mathrm{IDCG@K}(q_t)}.
\]

If $\mathrm{IDCG@K}(q_t)=0$, we set $\mathrm{nDCG@K}(q_t)=0$.

\subsection{Controller Parameters}
\label{sec:12-appendix:controller-parameters}

The complete shared controller configuration is:

\begin{itemize}
\item KMeans/MiniBatchKMeans arms: 32 fixed LinUCB arms;
\item candidate arms per query: 3 cluster-local routes;
\item context projection dimension: 64;
\item LinUCB exploration: $\alpha=1.0$, decay 0.01, minimum 0.3;
\item prequential epochs: 3 unless otherwise stated;
\item feedback trust: default $\tau=0.75$ for noisy feedback updates;
\item full-route candidate depths: dense/BM25/cluster = 100/100/100;
\item lite-route depths: dense/BM25 = 20/20;
\item dense safety floor: 5 full-route chunks and 2 lite-route chunks;
\item route fusion: weighted reciprocal-rank fusion with $k=60$;
\item full-route weights: dense 2.0, BM25 0.8, cluster 0.8;
\item lite-route weights: dense 0.8, BM25 0.5, cluster 2.0;
\item confidence thresholds: high 0.65 and mid 0.35;
\item drift threshold: 1.0;
\item token-budget grid: $r \in \{0.85,0.88,0.90,0.92,0.95,0.98\}$ and $m \in \{4,\ldots,8\}$.
\end{itemize}

The notation \nolinkurl{token_budget_r0.85_m4} means that each query keeps a safe prefix of at least four chunks and admits additional chunks only while the final context remains within 85\% of the original dense top-10 token budget. The policy is selected on calibration queries and frozen before held-out test evaluation.

\subsection{Reporting Guardrails}
\label{sec:12-appendix:reporting-guardrails}

The following rules apply when migrating the draft into a submission template:

\begin{itemize}
\item report query-level $\mathrm{Hit@10}$ as the primary retrieval headline;
\item report $\mathrm{EvidenceRecall@10}$ separately for complete-evidence tasks;
\item use final retrieved context tokens for prompt-context efficiency claims;
\item label source candidate cost and dense invocation rate as retrieval-stage diagnostics;
\item describe multi-epoch prequential adaptation as simulated repeated interaction, not IID held-out generalization;
\item describe feedback as controlled simulation, not collected production feedback;
\item describe geometry diagnostics as support for a piecewise local-structure interpretation, not theorem-level manifold proof;
\item keep dense retrieval visible as a recall floor and fallback route.
\end{itemize}

\subsection{Result-Family Protocol Registry}
\label{sec:12-appendix:result-family-protocol-registry}

\begin{table}[pos=tbp]
\centering
\small
\caption{Dataset, protocol, and evidentiary-role registry}
\label{tab:s29}
\begin{tabularx}{\textwidth}{>{\raggedright\arraybackslash}X>{\raggedright\arraybackslash}X>{\raggedright\arraybackslash}X>{\raggedright\arraybackslash}X>{\raggedright\arraybackslash}X}
\toprule
Result family & Corpus and query scope & Route/feedback protocol & Budget and paired endpoint & Paper role and boundary \\
\midrule
LoTTE technology/search scale anchor & 100k--638k nested corpora; 596 canonical queries (417 in the original frozen split; 596 in OOF) & MiniLM, top-10, K=32, seeds 13/17/19; frozen scale-route artifacts & Original 30/70 calibration/test plus normalized five-fold OOF; paired bootstrap, McNemar, and 1pp NI & Main scale evidence; original and OOF results are reported separately \\
LoTTE science/search diagnostic & 20k/q200; 200 sampled queries & Historical MiniLM fixed-split route study & Fixed 30/70 calibration/test & Legacy cross-domain ranking and budget diagnostic; not pooled with OOF rows \\
LoTTE science/search common rows & 100k, 200k, 400k; 596 queries per scale & MiniLM, top-10, K=32, seeds 13/17/19; 8 prequential epochs; final-fused simulated feedback with no-feedback control & Five-fold OOF zero-drop budget selection, Dense fallback, paired bootstrap, McNemar, 1pp NI, and recovery replay & Cross-domain evidence; 400k is a weak scale boundary (1/5 eligible folds) \\
LoTTE recreation/search and writing/search expansion & 100,714/924 and 100,696/1,071 corpus/query pairs & MiniLM, top-10, K=32, seeds 13/17/19; 8 prequential epochs; trust-weighted and no-feedback controls & Five-fold OOF zero-drop budget selection, Dense fallback, paired bootstrap, McNemar, 1pp NI, geometry diagnostics, and post-failure recovery & Preregistered 100k external-validity test; writing is a useful no-feedback frontier, recreation is a weaker boundary, and trust-weighted calibration falls back in both; no pooling \\
Frozen-policy unseen-query audit & LoTTE technology/search and science/search 100k; 596 queries each & Five history/test folds; seeds 13/17/19; eight history epochs; policy and feedback state frozen on the held-out fold & Query-paired bootstrap and McNemar on retrieval metrics; no held-out feedback update or budget endpoint & Boundary evidence: full rescue surface transfers, but learned feedback does not beat matched static/cold full routing and frozen gating is unsafe \\
PubMedQA, CovidQA-RAG, eManual deduplicated & Native full: 4,348/1,000; 32,392/1,726; 1,729/130 corpus/query pairs & MiniLM, top-10, K=32, seeds 13/17/19; 8 prequential epochs; final-fused trust/no-feedback controls & Five-fold OOF zero-drop budget selection and paired statistics & Transfer and corrected-boundary rows; not a shared scale curve \\
Banking77 and CUAD & Banking77 native 10,003/3,080; CUAD GT-anchored 10k/79 evaluated & Historical route-learning proxy and sparse-GT smoke protocols & No common OOF token endpoint & Mechanism and boundary evidence only; excluded from pooled evidence-retrieval conclusions \\
\bottomrule
\end{tabularx}
\end{table}

Table~\ref{tab:s29} is the canonical registry for the matched protocol definitions used in the cross-dataset comparisons.

The registry describes the paper-facing result families rather than retroactively equating their tasks. Detailed artifact hashes, historical protocol labels, and all route parameters remain in the reproducibility audit under \nolinkurl{paper/experiments/}.

\section{Preregistered LoTTE Domain Expansion}
\label{sec:12-appendix:preregistered-lotte-domain-expansion}

The domain-expansion protocol selected recreation/search and writing/search before downloading or inspecting their outcomes, then retained both after the preregistered lexicality ordering was contradicted. Recreation/search has lower query-positive token coverage than writing/search (0.6755 versus 0.7478) and lower maximum Jaccard overlap (0.0787 versus 0.1003). The corresponding recreation-minus-writing bootstrap intervals are [-0.0901, -0.0543] and [-0.0268, -0.0165]. The original lexicality premise is therefore not used to explain the budget result. Table~\ref{tab:s30} reports the matched route and budget outcomes.

\begin{table}[pos=tbp]
\centering
\small
\caption{Preregistered LoTTE 100k domain expansion. Full-route Hit precedes independent five-fold context-budget calibration; trust rows with zero eligible folds use Dense top-10 fallback}
\label{tab:s30}
\begin{tabularx}{\textwidth}{>{\raggedright\arraybackslash}X>{\raggedleft\arraybackslash}X>{\raggedleft\arraybackslash}X>{\raggedleft\arraybackslash}X>{\raggedleft\arraybackslash}X>{\raggedright\arraybackslash}X>{\raggedright\arraybackslash}X>{\raggedleft\arraybackslash}X>{\raggedright\arraybackslash}X}
\toprule
Domain / control & Lexical coverage & Dense Hit & Full-route Hit & Nearest-cluster Hit@3 & Eligible folds & Budget Hit delta & Token saving & Strict NI seeds \\
\midrule
recreation / trust & 0.6755 & 0.8496 & 0.8597 & 0.8366 & 0/5 & +0.00 pp & 0.00\% & n/a (fallback) \\
recreation / no feedback & 0.6755 & 0.8496 & 0.8546 & 0.8366 & 4/5 & -0.76 pp & 5.42\% & 0/3 \\
writing / trust & 0.7478 & 0.8739 & 0.8842 & 0.8655 & 0/5 & +0.00 pp & 0.00\% & n/a (fallback) \\
writing / no feedback & 0.7478 & 0.8739 & 0.8842 & 0.8655 & 5/5 & +0.12 pp & 10.09\% & 2/3 \\
\bottomrule
\end{tabularx}
\end{table}

Both domains contain usable cluster-local route signal, but their calibrated frontiers and strict seed-level stability differ. The result supports bounded external validity and domain-specific calibration. It does not establish that geometry or trust-weighted feedback directly predicts compression safety, and it does not convert Dense fallback into a positive token-saving result.

The post-failure replay remains supporting evidence only. Across seed-level compression-only failures, arm boost recovers 12/40 events on recreation/search and 28/29 on writing/search. Because the retry observes the same query's simulated failure feedback, it is not first-pass unseen-query generalization.
